% !TeX root = ../main.tex

\chapter{Operational Safety e Agent Guardrails}
\label{cap:safety_guardrails}

L'introduzione di agenti LLM espone il sistema a rischi (allucinazioni, latenze, comportamenti non deterministici): il progetto adotta più livelli di guardrail per garantire azioni verificabili, confinate e sicure.

\section{Least Privilege e Protezione Anti-Loop}
\label{sec:least_privilege}

L'IA agisce da \emph{assistente operativo}, mai con controllo illimitato. Il server MCP è l'unico \textit{chokepoint} verso l'infrastruttura e applica:
\begin{itemize}
    \item \textbf{Accesso read-only alla telemetria:} solo \texttt{query\_api} pre-compilate in Python, senza facoltà di comporre query \textit{Flux}/\textit{SQL} arbitrarie (rischio injection annullato).
    \item \textbf{Divieto di azioni distruttive:} nessun tool per \texttt{DROP} di bucket, \texttt{DELETE} di ordini storici o terminazione di processi core; il codice in \texttt{drone\_mcp\_layer.py} omette intenzionalmente queste implementazioni.
\end{itemize}

Per evitare loop di ragionamento infiniti, l'\texttt{HealthAgent} ha un contatore di iterazioni per turno (limite 3) e un controllo di formato sulla risposta: due output consecutivi non conformi terminano il processo. È inoltre presente un \textit{kill-switch} contro lo scollegamento del backend: se le chiamate al server MCP falliscono oltre una soglia, l'attività degli agenti viene congelata e ripresa automaticamente al ritorno online del backend, evitando di consumare token su contesti vuoti.

\section{Idempotenza delle Azioni}
\label{sec:idempotenza_azioni}

L'idempotenza è una proprietà di sicurezza fondamentale per evitare fallimenti catastrofici dovuti a loop o allucinazioni. Il meccanismo di scaling è l'esempio canonico: \texttt{scale\_drone\_deployment} invoca \texttt{patch\_namespaced\_deployment\_scale} passando il \emph{numero di repliche desiderate} (stato desiderato dichiarativo). Se l'LLM ripete il comando in loop, Kubernetes verifica che il Deployment ha già raggiunto la quota richiesta e ignora le chiamate successive. Una progettazione \emph{incrementale} (aggiungere $N$ droni per chiamata) avrebbe invece, in caso di allucinazione, saturato il cluster in pochi secondi, congestionato la rete MQTT e violato i tetti di spesa.

\section{Human-In-The-Loop ed Escalation}
\label{sec:hitl_escalation}

\textbf{La regola architetturale dei 6 droni.} La costante \texttt{POLICY\_LIMITS} in \texttt{drone\_mcp\_layer.py} fissa a 6 il numero massimo di droni scalabili in autonomia (\texttt{max\_drones\_auto}):
\begin{itemize}
    \item $\le 6$: l'\texttt{HealthAgent} invoca direttamente \texttt{scale\_drone\_deployment}.
    \item $> 6$: il layer MCP blocca il comando e l'agente è obbligato a invocare \texttt{request\_human\_approval}.
\end{itemize}

\textbf{Il portale di approvazione.} Le richieste sono gestite da un servizio Flask indipendente (\texttt{human\_approval\_manager.py}, porta 5002) che espone una dashboard dedicata agli operatori. Quando l'agente richiede un'escalation, il server MCP persiste la richiesta in \texttt{pending\_approvals.jsonl} e la dashboard renderizza una scheda con \texttt{request\_id}, azione, payload JSON (es. \texttt{\{"replicas": 8\}}) e, soprattutto, la giustificazione in linguaggio naturale fornita dall'LLM (\texttt{reason}). L'operatore dispone di un diritto di veto assoluto: \textbf{Approve} convalida la transizione ed esegue istantaneamente il patch su Kubernetes con \texttt{force=True}, scavalcando i blocchi preventivi; \textbf{Deny} contrassegna la richiesta come \texttt{rejected} e al ciclo successivo l'\texttt{HealthAgent} mantiene lo stato della flotta inalterato.
